\subsection{The twenty-five legal corrections of the first five-primary gate}

\begin{result}{E.202}{Proved}
{An exact affine system for the first two original factorial support primes}
Fix a legal first pair $q=q_0,r=r_0$ in the minimal-five-shell
subfamily of E.199, so $v_5(q-1)=v_5(r-1)=1$.
There is a subfamily of the unchanged recursion E.102, indexed by
$\mathbb F_5^2$, with the following properties.  Its global abelian
cochains modulo five are fixed.  Every parameter occurs by changing
only the fifth-root-of-unity factors in the two old local roots at
stage one, followed by legal choices in the original recursion.
The first two coefficient-five obstruction values depend only on
that parameter, not on the later prime choices in this subfamily.
Moreover these twenty-five parameters exhaust the possible first-gate
values of all original-recursion continuations with the fixed $q,r$.

Here is their exact system.  Fix the local tame generators and
arithmetic Frobenius lifts at $q,r$, and use the resulting common
Tate coordinate.  Put $m=24\in\mathbb F_5^\times$.
There are unique classes
\[
 z_q,z_r\in H^1(G_{\mathbb Q,\{5,q,r,\infty\}},\mathbb F_5(2))
\]
whose tame residue vectors are $(1,0)$ and $(0,1)$.
Choose any one branch as reference.  Write $D_q,D_r$ for its two
E.199 mismatch values modulo five and $\alpha=a(F_q)$,
$\beta=b(F_r)$ modulo five.  Then the first two local gates are
simultaneously satisfied by some original-recursion continuation
with the fixed $q,r$ if and only if
\[
 \begin{pmatrix}
 z_q(F_q)-\alpha/m&z_r(F_q)\\
 z_q(F_r)&z_r(F_r)-\beta/m
 \end{pmatrix}
 \begin{pmatrix}u\\v\end{pmatrix}
 =-\begin{pmatrix}D_q\\D_r\end{pmatrix}
 \quad\text{over }\mathbb F_5.
 \tag{E.202.1}
\]
If consistent, this gives exactly $5^{2-\operatorname{rank}M}$
root-choice parameters, where $M$ is the displayed matrix.
Neither consistency nor the later local and strict $c_0$ conditions
are asserted.  In particular this result does not construct a
class-three survivor or an all-depth branch.
\end{result}

\begin{proof}
\emph{The exact finite cohomology space.}
For $S=\{5,\infty\}$ the fixed-ramification calculations underlying
A.20 and A.22 give
\[
 H^1(G_{\mathbb Q,S},\mathbb Z_5(2))=0,
 \qquad H^2(G_{\mathbb Q,S},\mathbb Z_5(2))=0.
 \tag{E.202.2}
\]
The second equality is the five-primary part of the weight-two
arithmetic group $K_2(\mathbb Z[1/5])$, which is zero; localization
adds no five-primary part to $K_2(\mathbb Z)=\mathbb Z/2$.
The coefficient-five cohomological dimension is two, and the
finite-coefficient cohomology groups for this fixed $S$ are finite.
Thus the coefficient inverse limit has zero degree-three group:
the possible $\varprojlim^1$ of degree-two finite groups vanishes.
The multiplication-by-five coefficient sequence applied to
E.202.2 consequently gives
\[
 H^1(G_{\mathbb Q,S},\mathbb F_5(2))=
 H^2(G_{\mathbb Q,S},\mathbb F_5(2))=0.
\]
The finite-coefficient purity sequence, used with the same tame
normalization as A.22 and B.22, therefore makes the residue map
\[
 H^1(G_{\mathbb Q,S\cup\{q,r\}},\mathbb F_5(2))
 \ \xrightarrow{\sim}\
 H^0(\mathbb F_q,\mathbb F_5(1))\oplus
 H^0(\mathbb F_r,\mathbb F_5(1))=\mathbb F_5^2.
 \tag{E.202.3}
\]
This proves existence and uniqueness of $z_q,z_r$.
At these two places $\chi_5^2=1$ modulo five.  The value on a
fixed Frobenius of a cohomology class is consequently independent
of its cocycle representative.  The specified tame generators
fix the two residue normalizations.

\emph{Legal root choices in the original recursion.}
The first stage has already produced
$a_1=3q^{24}$ and $b_1=-2r^{24}$.
At stage one the exponent is $n_1=120$ and the next exponent is
$n_2=720$.  The local equation is
\[
                  a_1X^{120}+b_1Y^{120}=1.
 \tag{E.202.4}
\]
At $q$ choose any old unit solution $(\xi_q,\eta_q)$ and
replace $\eta_q$ by $\eta_q\zeta_q^i$, where $\zeta_q$ is a
primitive fifth root of unity in $\mathbb Q_q$ and
$i\in\mathbb F_5$.  At $r$ independently replace $\xi_r$ by
$\xi_r\zeta_r^j$, leaving $\eta_r$ fixed.  Since $5\mid120$,
these remain exact solutions of E.202.4.
The primes $q,r$ are distinct from $2,3,5$, and hence do not
divide $720$.  No condition requiring the old roots to be one
modulo a prime power dividing $720$ is changed.  All other old
roots, including the exact five-adic shells of E.199, can be kept.
The original CRT and complete-splitting argument of E.102 then
chooses $q_1$ and $r_1$ in all twenty-five cases.  In particular
its current-prime radical and congruence at $q_1$ are retained;
there is no replacement by a different prime-selection theorem.

After this step, the new coefficient sum is one to the prescribed
positive precision at both $q$ and $r$.  At every subsequent
stage choose both old roots congruent to one modulo the relevant
old prime.  This is always possible.  At $q$ the coefficient of
$Y^{n_j}$ is a unit, the other coefficient is divisible by $q$,
and $q\nmid n_j$ while this simple-root argument is needed;
setting $X=1$ and applying Hensel gives $Y\equiv1\pmod q$.
If eventually $q\mid n_j$, the original stronger error bound and
Newton estimate in E.102 give the roots congruent to one modulo
$q^{v_q(n_{j+1})}$, which in particular gives the stated
congruence.  The same argument at $r$, using the unit coefficient
of $X^{n_j}$, applies there.  All other old-prime conditions and
the minimal-five-shell choices can still be imposed separately.
Thus every parameter admits a complete legal recursion.

\emph{What those roots change.}
Choose global integral Kummer cochains so that their reductions
modulo five are the fixed cochains of $3q^{24}$ and $-2r^{24}$.
This is possible simultaneously on every branch: all later
exponents are divisible by five, and the sums of the integral
cochains multiplied by the factorial exponents converge uniformly.
Let $h,h'$ be the integral class-two cochains for two branches.
At coefficient five their class is unique, because
$H^1(G_{\mathbb Q},\mathbb Z_5(2))=0$ by A.20.
Their difference modulo five is therefore a well-defined
cohomology class $z$ with coefficients $\mathbb F_5(2)$;
its coboundary is zero because the two abelian cochains agree
modulo five.  Global central coboundaries do not affect any of
the displayed local Frobenius values.

The class-two tame relation at an $A$-support prime, in the
notation of B.102, is
\[
                  h(T)=n_j u,\qquad b(F)=(p-1)u;
\]
at a $B$-support prime it is
\[
                  h(T)=-n_j v,\qquad a(F)=(p-1)v.
 \tag{E.202.5}
\]
The divided primitives are integral because the complementary
ordinary coordinate is a principal unit.  Hence every new
support prime, whose exponent is divisible by five, has zero
$h(T)$ modulo five.  At zero-support primes outside five the
same tame relation gives $(p-1)h(T)=0$ integrally, and so
$h(T)=0$.  The fixed primes two and three likewise introduce
no residue for the difference: their finite residue spaces
$H^0(\mathbb F_p,\mathbb F_5(1))$ vanish.  It follows that
$z$ is unramified outside $\{5,q,r,\infty\}$.

Fix compatible roots of unity in the two local embeddings.
The Kummer cochain of $\zeta_q$ has Frobenius value
$(q-1)/5$, up to an integral coboundary.  Multiplication by
$120$ therefore changes the divided primitive $u_q$ by
$24$ modulo five.  Any principal-unit factor has an integral
Kummer primitive, whose multiple by $120$ contributes zero
modulo five.  The CRT precision is at least one at $q$,
so the ratio of the two selected $r_1$ residues is exactly the
specified fifth root of unity modulo $q$; all further factors
are principal units.  Thus the $i$-choice changes $u_q$ by
$24i$ modulo five.  Independently the $j$-choice changes
$v_r$ by $24j$.  The same argument shows that all subsequent
old roots congruent to one at $q,r$ make zero contribution to
these changes modulo five.  This conclusion also covers the
finite precision in every CRT choice.

By E.202.5 the residue vector of $z$ is
\[
                 (24\,\Delta u_q,-24\,\Delta v_r).
 \tag{E.202.6}
\]
Both scalar factors, and the factors $24$ just obtained from
the root choices, are units modulo five.  The twenty-five root
choices therefore give every vector in $\mathbb F_5^2$ exactly
once.  By E.202.3, writing that vector as $(u,v)$ gives
$z=u z_q+v z_r$.  In particular no other later prime choice
inside the stated subfamily can alter this class.

\emph{The central equations.}
At $q$ the E.199 mismatch is $h(F_q)-q\alpha u_q$, and at
$r$ it is $h(F_r)+\beta v_r$.  The fixed abelian cochains
make $\alpha,\beta$ common to all branches modulo five.
Since $q\equiv r\equiv1\pmod5$, E.202.6 gives the changes
\[
\begin{aligned}
 \Delta D_q&=z(F_q)-\alpha u/24,\\
 \Delta D_r&=z(F_r)-\beta v/24.
\end{aligned}
\]
Substitute $z=u z_q+v z_r$ to obtain E.202.1.
For an arbitrary original-recursion continuation with the same
$q,r$, the identical inertia argument puts its difference from the
reference in E.202.3.  Its primitive changes are still E.202.6.
Thus its first-gate vector is already one of those attained by the
constructed twenty-five choices, proving the asserted exhaustion.
An alternative Frobenius lift changes its row by the same tame
residue correction in the cochain and the abelian term, so the
criterion retains its meaning in the stated fixed frames.
Elementary linear algebra gives the count of successful
parameters.  The equations at later indices divisible by five
and the strict convergence condition of E.199 remain additional
requirements, and no inference concerning them is made here.
\end{proof}

\paragraph{Explicit Kummer coordinates for the matrix.}
The entries of E.202.1 admit the following finite-field computation,
with its ramified Frobenius lifts fixed explicitly.  Let
$K=\mathbb Q(\zeta)$, where $\zeta$ is a primitive fifth root,
and write $\sigma_a(\zeta)=\zeta^a$ for $1\le a\le4$.
For each $p\in\{q,r\}$ choose the prime $\mathfrak p$ of $K$
specified by the local embedding, and a generator $\pi_p$ of it.
Put
\[
 c_p=\prod_{a=1}^4\sigma_a(\pi_p)^a,
 \qquad U_p=c_p/p.
 \tag{E.202.7}
\]
Then $U_p$ is a unit at $\mathfrak p$.
If $\omega_p\in\mathbb F_p^\times$ is the reduction of $\zeta$,
define the quintic residue logarithm of a $\mathfrak p$-unit by
\[
       \overline t^{\,(p-1)/5}=\omega_p^{\,\ell_p(t)},
       \qquad \ell_p(t)\in\mathbb F_5.
 \tag{E.202.8}
\]
Choose $F_p$ to fix compatible fifth-power roots of the rational
uniformizer $p$ in the tame extension.  The matrix is then exactly
\[
 M=\begin{pmatrix}
       \ell_q(3U_q)&\ell_q(c_r)\\
       \ell_r(c_q)&\ell_r(-2U_r)
    \end{pmatrix}.
 \tag{E.202.9}
\]
Thus no unspecified ramified Frobenius value remains in its
calculation.  This formula supplies the matrix, not the reference
constant vector $(D_q,D_r)$ or a proof that it is in the image.

Here are the arithmetic and descent details.  The field $K$ has
class number one: its discriminant has absolute value $125$, and
Minkowski's bound
\citep[Theorem4.3, p.~70]{MilneANT2020} is
\[
       \frac{4!}{4^4}\left(\frac4\pi\right)^2\sqrt{125}<2.
\]
Every ideal class therefore has an integral representative of norm
one.  The chosen prime generators have norm $p$, since $K$ is
CM and the norm of a generator of a degree-one prime has absolute
value $p$.  In particular $\prod_a\sigma_a(\pi_p)=p$.
The valuations of $c_p$ at the four conjugate primes are
$1,2,3,4$ in the corresponding order, so $U_p$ is indeed a unit
at the chosen prime.

In $K^\times/K^{\times5}$ one has
\[
              \sigma_b[c_p]=b^{-1}[c_p].
 \tag{E.202.10}
\]
This follows by reindexing $a$ by $ba$; replacing the exponent
$b^{-1}a$ by its representative from $1$ to $4$ changes only a
fifth power.  The same weighted product annihilates the unit group
modulo fifth powers.  For detail, on units modulo roots of unity
complex conjugation acts trivially.  Indeed a unit divided by its
complex conjugate is an algebraic integer with all conjugates of
absolute value one, hence a root of unity: its powers have bounded
integral characteristic-polynomial coefficients, so two powers
coincide.  The only five-primary roots of unity in $K$ form
$\mu_5$, on which $\sigma_b$ acts as $b$.  The $b^{-1}$
eigenspace is odd under complex conjugation and differs from this
torsion character.  Since the group of order four is semisimple
in characteristic five, that eigenspace of the units is zero.
This also proves independence of the class $[c_p]$ from the chosen
prime generator.  The valuation at the unique prime above five
has trivial Galois character, so it contributes no such eigenspace
either.

By Kummer theory, $\kappa_5(c_p)\otimes\zeta$ is an invariant
class in $H^1(G_K,\mathbb F_5(2))$.  Inflation--restriction for
$\operatorname{Gal}(K/\mathbb Q)$, whose order is prime to five,
descends it uniquely to $H^1(G_{\mathbb Q},\mathbb F_5(2))$.
It is unramified outside $5,p$, and has residue one at the chosen
$p$ embedding.  By E.202.3 it is exactly $z_p$.
For the specified Frobenius, the Kummer cochain of $p$ is zero,
and the Kummer cochain of a unit has value E.202.8: if $w^5=t$,
then arithmetic Frobenius gives $w^{p-1}=t^{(p-1)/5}$ after
reduction.  Hence $z_p(F_p)=\ell_p(U_p)$ and the off-diagonal
values are $z_q(F_r)=\ell_r(c_q)$ and
$z_r(F_q)=\ell_q(c_r)$.
The fixed abelian classes are those of $3q^{24}$ and $-2r^{24}$;
therefore $\alpha=\ell_q(3)$ and $\beta=\ell_r(-2)$ in these
frames.  Finally $24^{-1}=-1$ in $\mathbb F_5$, giving E.202.9.
No numerical legal first prime pair or matrix rank is claimed.

\paragraph{The constant-vector class also has a finite description.}
The reference continuation in E.202.1 is unnecessary if only the
obstruction class in $\operatorname{coker}M$ is wanted.  Put
$S_0=\{2,3,5,q,r,\infty\}$ and retain the finite Kummer pair
\[
          a=\kappa_5(3q^{24}),\qquad
          b=\kappa_5(-2r^{24}).
 \tag{E.202.11}
\]
Its cup is zero already in
$H^2(G_{\mathbb Q,S_0},\mathbb F_5(2))$.
Indeed finite purity and the vanishing of the two fixed-$S$
groups in the proof identify this degree-two group with the
residue spaces $H^1(\mathbb F_p,\mathbb F_5(1))$ for
$p=2,3,q,r$.  The spaces at two and three vanish.  At $q$,
the original current-prime congruence for $r_0$ gives
$-2r^{24}\equiv1\pmod q$, so $b(F_q)=0$ modulo five.
At $r$, the original splitting in
$\mathbb Q(\mu_{120},(3q^{24})^{1/120})$ gives
$a(F_r)=0$ modulo five.  The tame cup residues are therefore
zero at both remaining places.  The degree-three term at $S$
is zero by its coefficient-five cohomological dimension, so
there is no further term in this purity argument.

Choose any finite class-two cochain $h_0$ above E.202.11 over
$G_{\mathbb Q,S_0}$, in the same Hall convention as B.102.
The cup vanishing gives such a cochain.  Define
\[
 d_q=h_0(F_q)-\frac{\alpha}{24}h_0(T_q),\qquad
 d_r=h_0(F_r)-\frac{\beta}{24}h_0(T_r).
 \tag{E.202.12}
\]
Then $(d_q,d_r)$ represents exactly the constant-vector class
of E.202.1 in $\operatorname{coker}M$.
To verify this, the reduction of any integral continuation is
also a finite class-two cochain over $G_{\mathbb Q,S_0}$:
its abelian inertia vanishes at all later support primes modulo
five, and E.202.5 makes its central inertia vanish there too.
The fixed arithmetic action is unramified outside five.
Its difference from $h_0$ is consequently in E.202.3.
For an integral continuation, E.202.5 identifies E.202.12 with
its $D_q,D_r$.  Changing $h_0$ by $u z_q+v z_r$ changes
E.202.12 by precisely $M(u,v)$.  All such classes are realized
by the twenty-five legal root choices already proved.

Thus the first two tests, over all continuations of a fixed legal
$q_0,r_0$, are a finite arithmetic problem depending only on
that first prime pair.  No later choice can repair an inconsistent
system.  The formulas do not evaluate that system, produce a
certified legal numerical first prime pair, or settle any later
local residue or the strict $c_0$ condition.
